\section{Award-Layer Triage}
\label{sec:award_screen}

The award layer observes outcomes before it observes conduct. A
losing bid may be sincere, poorly priced, capacity constrained,
exploratory, strategic, or coordinated. The screen therefore
uses what the award layer can measure at the triage stage: which
firms appeared, which firms won, and which firms kept appearing
without ever winning. The object is a ranking for forensic
priority, not proof of conduct or a classification of cartel
membership.

This section defines that ranking and the empirical implications
that discipline it. Subsection~\ref{sec:award_primitive}
introduces persistent zero-win participation as the award-layer
primitive. Subsection~\ref{sec:award_rule} explains why the
continuous ranking can be converted into an operational queue or
binary frequent-loser flag. Subsection
\ref{sec:testable_implications} states the tests that the rest
of the paper uses to evaluate whether the ranking is informative
for enforcement.

% CO-AUTHOR EDIT [v21 conceptual framing]: decision-theoretic frame.
% Drafted+adversarially verified via workflow; anchors bound to the real
% exposure macros (\valExpoWithinStratumAUC etc.) and temporal TP macros.
\subsection{A Decision-Theoretic Frame for Award-Layer Triage}
\label{sec:triage-model}

We model the enforcement agency, not the conduct it screens. The
question is how an agency that already observes a cheap,
award-layer signal at negligible cost should decide \emph{whom to
escalate} to a costly evidentiary stage under a fixed recovery
budget. Nothing here turns on a theory of how rings bid; the
bid-layer price patterns enter the empirics as \emph{scope}, not
as a damages or mechanism claim. We take the informativeness of
the cheap signal as a maintained primitive---the monotone
likelihood ratio property (MLRP) stated, not derived, in
\S\ref{app:framework_participation}.

\paragraph{Primitives.}
Index zero-win firms by $i$, each with a latent enforcement target
$\theta_i\in\{0,1\}$ ($\theta_i=1$ = cartel-adjacent co-bidding
position; small base rate $\pi=\Pr(\theta_i=1)$). The agency
observes at $\approx 0$ cost a \emph{cheap} award-layer signal
$s_i=\log(1+\text{tenders\_count}_i)$ (loss intensity); a
\emph{costly} bid-layer signal $b_i$ (bid-distribution moments) is
recoverable only at cost $c>0$ per firm. Both are confounded by
\emph{exposure} $e_i$, an award-observable opportunity measure: by
a coupon-collector logic, larger $e_i$ raises both the chance of
having co-bid with a sanctioned firm ($\theta$) and the loss count
($s$). Write $\rho(s,e)=\Pr(\theta=1\mid s,e)$ and
$\rho_e(e)=\Pr(\theta=1\mid e)$.

\paragraph{The agency's problem.}
With recovery budget $K$, the agency escalates a set $A$ of at
most $K$ firms to the costly stage using only pre-recovery
information $(s,e)$, maximizing expected target yield
\begin{equation}
\label{eq:agency-yield}
V(A)=\sum_{i\in A}\rho(s_i,e_i),\qquad |A|\le K .
\end{equation}
A cheap gate orders the queue; the costly stage adjudicates its
head.

\begin{proposition}[Cheap-signal index rule]
\label{prop:index}
The yield \eqref{eq:agency-yield} is maximized over $|A|\le K$ by
escalating the $K$ firms with the highest posterior
$\rho(s_i,e_i)$. Under MLRP of $s\mid(\theta,e)$, $\rho(s,e)$ is
nondecreasing in $s$ for each fixed $e$, so within any exposure
stratum the optimal rule is a cutoff on the award rank of $s$,
deployable as a per-stratum top-$k$ queue.
\end{proposition}

\noindent\emph{Intuition.} A deployability floor: a budgeted
modular yield is maximized by a greedy index, and MLRP turns that
index into a threshold on an observable rank
(\ref{app:triage-proofs}). It underwrites the operational queue of
\S\ref{sec:award_rule} and the sequential gatekeeper, which traces
a cost--recall frontier and shrinks the costly-recovery footprint
by $\valArchFootprintReduction$ at bounded recall cost.

\begin{proposition}[Incremental value of the cheap signal over exposure]
\label{prop:incremental}
Fix the budget so that in at least one exposure stratum the agency
escalates some but not all firms (a binding budget within an
informative stratum), and let the exposure-only rule break
within-stratum ties uninformatively about $\theta$. Then, under
the maintained MLRP,
\[
V\!\big(\text{top-}K\text{ by }\rho(s,e)\big)>
V\!\big(\text{top-}K\text{ by }\rho_e(e)\big)
\]
if and only if $\theta$ is not conditionally independent of $s$
given $e$---equivalently, the within-exposure-stratum rank
discrimination of $s$ for $\theta$ exceeds chance,
$\mathrm{AUC}(s,\theta\mid e)>\tfrac12$. If $\theta\perp s\mid e$
the two policies coincide.
\end{proposition}

\noindent\emph{Intuition and caveat.} The ``if'' is immediate:
under $\theta\perp s\mid e$, $\rho(s,e)=\rho_e(e)$ and the policies
are identical. The equivalence ``conditional dependence
$\Leftrightarrow\mathrm{AUC}>\tfrac12$'' is \emph{not}
distribution-free---two conditional laws can differ yet rank-tie
at one half. It is the maintained MLRP that closes the
biconditional, forcing first-order stochastic dominance whenever
the laws differ; strictness in $V$ further requires the
exposure-only tie-break to be $\theta$-uninformative
(\ref{app:triage-proofs}). We state the equivalence \emph{within}
the informative, budget-binding stratum, not globally.

This is the result the data must test. The empirical content is
the within-opportunity rank discrimination
$\valExpoWithinStratumAUC$, which clears the one-half threshold of
Proposition~\ref{prop:incremental}, and the positive increment of
$+\valExpoIncr$ AUC it adds over an exposure-only model (DeLong
$p\approx\valExpoIncrP$). Exposure alone already reaches
$\valExpoExposureOnlyAUC$---\emph{above} the unconditional
$\valAUCFLfirm$---so most of the raw discrimination is
opportunity, not conduct, and the unconditional figure overstates
the conduct-relevant signal. The signal that survives is the
within-stratum $\valExpoWithinStratumAUC$ and the $+\valExpoIncr$
it adds over exposure alone. The condition of
Proposition~\ref{prop:incremental} holds in the data precisely
because this within-stratum increment is positive.

\begin{remark}[Out-of-sample value of the parsimonious gate]
\label{rem:oos}
In sample, a joint score $\rho_J(s,b)$ using both signals weakly
dominates any sequential gate that conditions first on $s$: the
gate is a garbling of $(s,b)$, so by the data-processing
inequality it cannot raise the population yield. Out of sample,
scores are estimated with error. The bid score aggregates many
estimated moments and inherits their variance; the cheap signal is
a single, low-variance count. When the estimation noise loaded
onto $\rho_J$ is large relative to the in-sample informational
gain $b$ buys, the sequential gate's expected out-of-sample yield
can exceed the joint score's, reversing the in-sample order. A toy
Gaussian (\ref{app:triage-proofs}) makes the bias--variance
crossing precise---and its limit honest: the reversal exists only
when $b$'s in-sample gain is smaller than the joint score's
maximal net noise penalty. We therefore state this as a Remark,
not a Proposition. The empirical pattern is consistent with the
high-relative-noise regime: under a temporal holdout, in the
budget range where the gate is designed to bind, the sequential
gate recovers at least as many targets as the joint score
($\valArchTHTPFiveHundSeqK$ versus $\valArchTHTPFiveHundJoint$ at
top-$500$; $\valArchTHTPThouSeqK$ versus $\valArchTHTPThouJoint$ at
top-$1{,}000$), while recovering microdata only for the gated
firms. We read this as consistent with the crossing, not as a
verification of the toy model, and note the advantage is confined
to the gating budget range.
\end{remark}

\begin{corollary}[Separation of attention from liability]
\label{cor:separation}
Let $D$ denote the direct defendants (the legal principals). If
role allocation makes the cheap signal approximately independent
of principal status, $s\perp\mathbf{1}[i\in D]$, then the triage
index is informative for \emph{escalation} ($\theta$) while
remaining near-random for \emph{liability}:
$\mathrm{AUC}(s,\mathbf{1}[i\in D])\approx\tfrac12$.
\end{corollary}

\noindent\emph{Intuition (the institutional hook).} This is the
property a well-designed enforcement institution wants: a triage
index that directs scarce evidentiary attention to the loser-side
target without being re-usable as a liability instrument, so it
cannot be weaponized as proof of guilt. Empirically the index is
near-random against direct defendants,
$\mathrm{AUC}=\valAUCdirectStd$---by construction, not by failure.
The premise $s\perp\mathbf{1}[i\in D]$ is an assumption about role
allocation; the observed near-random value is evidence consistent
with it.

\begin{table}[htbp]
\centering
\caption{From model primitives to empirical anchors.}
\label{tab:model-empirics}
\footnotesize
\begin{adjustbox}{max width=\textwidth,center}
\begin{threeparttable}
\begin{tabular}{p{3.3cm}p{5.3cm}p{3.6cm}}
\toprule
Model object & Empirical counterpart & Result / anchor \\
\midrule
Cheap signal $s_i$ & Loss intensity $\log(1+\text{tenders\_count})$, observed at $\approx 0$ cost & Prop.~\ref{prop:index} \\
Costly signal $b_i$, cost $c$ & Bid-distribution moments, recovered only on escalation & Rem.~\ref{rem:oos} \\
Exposure $e_i$ & Award-observable opportunity stratum & Prop.~\ref{prop:incremental} \\
Budget $K$, yield $V$ & Greedy top-$k$ queue; cost--recall frontier & Footprint $-\valArchFootprintReduction$ \\
$\mathrm{AUC}(s,\theta\mid e)>\tfrac12$ & Within-opportunity $\valExpoWithinStratumAUC$; $+\valExpoIncr$ over exposure-only ($p\approx\valExpoIncrP$); exposure alone $\valExpoExposureOnlyAUC$, residual to $\valAUCFLfirm$ & Prop.~\ref{prop:incremental} \\
Gate vs.\ joint, out of sample & Temporal holdout, gating range: $\valArchTHTPFiveHundSeqK$/$\valArchTHTPFiveHundJoint$ (top-500), $\valArchTHTPThouSeqK$/$\valArchTHTPThouJoint$ (top-1000) & Rem.~\ref{rem:oos} \\
$s\perp\mathbf{1}[i\in D]$ & Direct-defendant ranking, near-random & $\mathrm{AUC}=\valAUCdirectStd$ (Cor.~\ref{cor:separation}) \\
\bottomrule
\end{tabular}
\end{threeparttable}
\end{adjustbox}
\end{table}

\subsection{Persistent Zero-Win Participation}
\label{sec:award_primitive}

Let $T_i$ denote the number of BEC tender-items in which firm
$i$ participates, and let $W_i$ denote the number it wins. The
award-layer screen operates inside the zero-win stratum,
$\{i:W_i=0\}$, which contains $\valAlwaysLosers$ firms. Within
that stratum, the continuous score is
\[
  s_i=\log(1+T_i).
\]
The log transformation stabilizes the scale but does not change
the ordering. The economic object is the rank: among firms the
award layer observes only as losers, which ones appear
persistently enough to warrant bid-layer follow-up?

The maintained condition is monotonic rather than classificatory.
If cartel-deployed losing participants are present, they should
appear more persistently in the zero-win stratum than ordinary
losing firms, because such a role is useful only through repeated
deployment in cartel-targeted tenders. The condition is
maintained, not derived. \ref{app:framework_submission}
formalizes it as a monotone-likelihood-ratio ranking result. The
main text uses it more modestly: persistent zero-win participation
is a cheap statistic that can rank loser-side exposure under a
transparent behavioral assumption and direct attention to the bid
layer where conduct can be evaluated. It does not assign a firm
type or cartel membership.

The zero-win restriction is therefore a scope restriction. Many
firms never win for ordinary competitive reasons: poor fit, high costs,
capacity constraints, market exploration, subcontracting
positioning, or simple bad luck. The screen's content comes from
persistence inside that heterogeneous loser-side set. If the
ranking has enforcement value, it should concentrate
adjudication-anchored loser-side adjacency beyond what generic
participation volume, product-market overlap, and opportunity to
meet CADE anchors would mechanically produce.

\subsection{From Ranking to an Operational Queue}
\label{sec:award_rule}

The continuous score is the central object because agencies can
translate a rank into different investigative protocols. A team
with fixed capacity could inspect the top $k$ firms. A
monitoring unit could refresh the rank each quarter and examine
bunching below old cutoffs. A case team could use the score to
request LANCES bid files for a survivor pool before applying
bid-distribution forensics. These implementations differ, but
all preserve the same sequencing logic: award records first,
costly bid recovery second.

For an auditable binary implementation, we also report a
frequent-loser flag. The rule marks zero-win firms whose
participation count exceeds the median plus $1.5$ times the
interquartile range of the always-loser participation
distribution. In BEC, the statistical cutoff is
$\valThresholdStat$, implemented as $T_i\geq\valThreshold$, which
flags $\valFL$ firms, or $\valFLrateAL$ of the always-loser
stratum. The cutoff is fixed from the participation distribution
alone, before any CADE label enters the analysis.

The cutoff is an auditable administrative implementation of the
rank, not a structural or legal threshold. The paper does not
rest on the number $14$; it rests on the continuous loser-side
participation ranking. Two analysts applying the same rule to
the same award records obtain the same first-stage list. Above
the cutoff does not mean cartel member; below it does not mean
clean. Because a static threshold can be gamed, the enforcement
architecture treats the binary flag as one implementation among
several: agencies can also use continuous ranks, top-$k$ queues,
refreshed thresholds, or budget-calibrated survivor pools. The
strategic-adaptation analysis later returns to this point.

\subsection{Testable Implications}
\label{sec:testable_implications}

The screen is useful if it survives the threats implied by its
limited information. Table~\ref{tab:testable_implications_submission}
organizes the validation map. The first implication is the core
validation claim: persistent zero-win participation should rank
adjudication-anchored loser-side cobidders within the always-loser
stratum. The main threat is mechanical exposure: frequent
participants, and firms operating in the same procurement
environments as CADE defendants, have more opportunities to
appear near them even under ordinary competition.

The second implication is a scope check. A loser-side statistic
should perform worse against direct CADE defendants than against
always-loser cobidders, because it ranks a different role in the
procurement environment. The third implication asks whether the
cobidder target has economic content: among frequent losers,
cobidders should differ from other frequent losers in directions
consistent with loser-side exposure, proximity to defendants,
market concentration, and bid-level losing behavior. These
profiles give the validation target economic content; they do
not prove conduct or assign legal status.

The final implications evaluate the enforcement architecture.
If award records and bid records carry different information,
the award-layer score should add non-redundant signal to
bid-distribution forensics. If the screen has institutional
value, using it before bid recovery should reduce the
bid-microdata footprint while preserving much of the
adjudication-anchored recall. And because firms can respond to
published administrative rules, the binary threshold should be
treated as vulnerable to strategic adaptation rather than as a
fixed rule immune to gaming.

\begin{table}[htbp]
\centering
\caption{Testable implications and validation threats}
\label{tab:testable_implications_submission}
\footnotesize
\begin{adjustbox}{max width=\textwidth,center}
\begin{threeparttable}
\begin{tabular}{p{3.2cm}p{4.4cm}p{4.3cm}p{3.2cm}}
\toprule
Implication & Prediction & Test & Main threat \\
\midrule
Loser-side ranking & Persistent zero-win participation ranks always-loser cobidders. & Conservative benchmark, timing audit, exposure-adjusted validation. & Participation volume, market overlap, and opportunity to meet CADE anchors. \\
Scope asymmetry & Performance is weaker for direct defendants than for cobidders. & Direct-defendant AUC and temporal direct-defendant check. & Misreading the screen as generic cartel-membership detection. \\
Economic profile & Cobidders differ from other frequent losers in loser-side exposure and bid patterns. & Firm-level and bid-level profile comparisons. & Product-market composition and ordinary high-volume losing. \\
Complementarity & Award-layer signal adds information to bid-layer forensics. & Same-target award, bid, and combined AUCs. & Full-observability benchmark confused with operational sequencing. \\
Gatekeeping & Award-first sequencing reduces bid-microdata recovery with limited recall loss. & Sequential survivor-pool exercise and cost-recall frontier. & Survivor-pool choice and arbitrary budget choice. \\
Strategic response & Static cutoffs are more fragile than ranks or refreshed queues. & Adaptation diagnostics. & Simulation is not an equilibrium model. \\
\bottomrule
\end{tabular}
\begin{tablenotes}
\footnotesize
\item \textit{Notes:} Predictions evaluate whether the award-layer
ranking is useful for evidence allocation. They do not define
legal membership or proof.
\end{tablenotes}
\end{threeparttable}
\end{adjustbox}
\end{table}

Together, these tests evaluate the screen at the stage for which
it is designed: the allocation of costly forensic attention
before liability evidence is assembled. The BEC cutoff is an
implementation detail; the portable object is the staged
architecture that uses cheap administrative records to discipline
where costly cartel forensics should begin.
